Our experiment pipeline consists of two stages: (1) self-supervised pretraining via next-frame prediction on synthetic videos, and (2) reinforcement learning in the game environment using the pretrained model as initialization.

\textit{Pretraining Dataset.} We construct a pretraining dataset of synthetic videos generated from a variation of Conway's Game of Life~\citep{gardner1970life}. To increase complexity beyond the standard binary cellular automaton, we run $K$ parallel Game of Life fields simultaneously and color-code each field in the rendered $96 \times 96$ RGB display. Videos are generated by simulating $T$ timesteps from random initial configurations, yielding a diverse corpus of trajectories with rich spatiotemporal dynamics. Using synthetic data provides a controlled setting with known ground-truth dynamics, enabling precise evaluation of whether the model has internalized the underlying transition rules.

\textit{Frame Crossing Game.} We define a reinforcement learning task called the \textit{Frame Crossing Game}. The agent must navigate from one side of the grid to the opposite side in minimum time. The background consists of evolving Game of Life cells: contact with any active cell terminates the episode, while reaching the opposite boundary yields a positive reward. A per-step penalty encourages speed. Crucially, since the colored cells evolve according to Game of Life dynamics, the agent must anticipate future states to plan a safe trajectory, directly linking the pretraining objective to downstream task performance.
\subsection{Environments}

Conway's Game of Life~\citep{gardner1970life} is a two-dimensional cellular automaton defined on an $H \times W$ binary grid.
Each cell is either \emph{alive} (1) or \emph{dead} (0), and the next state is determined by counting the cell's eight spatial neighbors:
\begin{itemize}[nosep]
    \item \textbf{Birth:} a dead cell becomes alive if it has exactly 3 live neighbors.
    \item \textbf{Survival:} a live cell remains alive if it has 2 or 3 live neighbors; otherwise it dies.
\end{itemize}
We denote this rule set B3/S23.

\paragraph{Parallel fields.}
To scale the complexity of the environment, we run $K$ independent Game of Life fields on the same $96 \times 96$ grid simultaneously, where $K \in \{1,2,3,4\}$.
Each field evolves according to B3/S23 independently; the state of each cell is therefore a binary vector of length $K$, and the cell is considered dangerous to the agent if \emph{any} field is alive at that position.
Increasing $K$ raises the density of active cells and makes collision-free navigation exponentially harder.

\paragraph{Properties of the standard rule.}
The standard Game of Life has two key properties.
First, the update rule is \emph{local}: a cell's next value depends only on its $3 \times 3$ neighborhood.
Second, the transition is \emph{deterministic} and \emph{first-order Markov}: the next frame is a pure function of the current frame, with no dependence on earlier history.
Together, these properties mean that the standard setting has no temporal component---no attention over time is needed to resolve the prediction task. A model with sufficient spatial receptive field (at least $3 \times 3$) can, in principle, predict the next frame perfectly from a single input frame.
The task is, however, non-trivial \emph{spatially}: the model must learn to aggregate information over each cell's local neighborhood and apply the correct birth/survival logic, making a convolutional architecture a natural fit.

\paragraph{Higher-order dynamics.}
To introduce a genuine temporal dimension, we extend the automaton with a \emph{history order} $M \geq 1$.
When $M > 1$, neighbor counts are summed across the current frame and the previous $M - 1$ frames, and the birth/survival thresholds are scaled accordingly (e.g., with $M = 2$, birth occurs at neighbor counts $\{3, 4\}$ instead of $\{3\}$).
Under this rule, the next frame depends on the last $M$ frames, breaking the first-order Markov property and requiring the model to attend across time.
We set the temporal context window $k = M$ so that the model receives exactly enough history to determine the transition.
This extension provides a controlled knob for studying when temporal attention contributes to prediction accuracy.
\subsection{Models}

The world model is a spatial-temporal transformer that predicts the next
grid frame given $k$ context frames.
%
The \textbf{input} is a tensor of shape $(B, k, H, W)$, where each cell
is a 5-bit token encoding the binary state of up to four GoL fields
(bits 0--3) and the player presence (bit 4), giving a vocabulary of
$V = 2^5 = 32$.
%
An \textbf{embedding} layer maps each token to a $d$-dimensional vector
via a learned lookup table and adds three learned positional embeddings
(row, column, and time), producing a representation
$\mathbf{h}^{(0)} \in \mathbb{R}^{B \times k \times H \times W \times d}$.
%
This representation is then passed through $L$ identical
\textbf{transformer blocks}.  Each block applies, in sequence:
(i)~a $3 \times 3$ depthwise convolution for spatial mixing (residual
connection + LayerNorm),
(ii)~pre-norm multi-head self-attention over the temporal axis
(independently per pixel, across the $k$ frames), and
(iii)~a pre-norm two-layer MLP (Linear $\to$ GELU $\to$ Linear).
All three sub-layers use residual connections.
%
After the final block, only the last time-step representation
$\mathbf{h}^{(L)}_{k} \in \mathbb{R}^{B \times H \times W \times d}$
is retained.
%
An \textbf{unembedding} layer---a LayerNorm followed by a linear
projection---maps each cell's $d$-dimensional feature to $V = 32$
logits.
%
The \textbf{output} is therefore a tensor of shape
$(B, H, W, 32)$ representing next-frame token predictions at every cell.
%
During RL, no separate policy head is added: for each of the five
actions $\{$stay, up, down, left, right$\}$, the log-probability that
the player token is present at the corresponding target cell---computed
as $\log \sum_{v \geq 16} p_v - \log \sum_{v} p_v$---serves as the
action logit for a 5-way categorical distribution.

\subsection{Learning}

Both stages use the AdamW optimizer~\citep{loshchilov2019decoupledweightdecayregularization} with weight decay $10^{-4}$ and gradient clipping at max norm $1.0$. Mixed-precision training (bfloat16) is enabled throughout.

\textit{Pretraining.}
We generate 500 videos of 300 frames each across initial densities $\{0.10, 0.15, \ldots, 0.50\}$, split 90/10 into training and validation sets.
The model is trained for 6 epochs with batch size 256 and a peak learning rate of $10^{-3}$, using a linear warmup over 200 steps followed by cosine decay to $1\%$ of the peak.

\textit{Reinforcement learning.}
We train for 10{,}000 episodes grouped into batches of 16 episodes per gradient update.
The learning rate is $10^{-4}$ when fine-tuning from pretrained weights and $3 \times 10^{-4}$ when training from scratch.
Episodes last at most 150 steps and terminate early on collision with any live cell or upon reaching the right boundary.
The reward signal comprises a per-step penalty $r_{\text{step}} = -0.1$ to encourage speed, a one-time milestone bonus of $3.0$ for each new rightmost column reached, and a win bonus of $50.0$ for reaching the far edge.

\paragraph{Pretraining Loss.}
The pretraining objective is the standard cross-entropy loss over the $V = 32$-class token vocabulary, averaged over all spatial positions:
\begin{equation}
    \mathcal{L}_{\text{pretrain}} = - \frac{1}{HW} \sum_{i=1}^{H} \sum_{j=1}^{W} \log p_\theta\!\bigl(y_{ij} \mid \mathbf{x}\bigr),
\end{equation}
where $\mathbf{x}$ is the $k$-frame input context and $y_{ij}$ is the ground-truth token at position $(i,j)$ in the next frame.

\paragraph{Reinforcement Learning Loss.}
The RL objective combines three terms:
\begin{equation}
    \mathcal{L}_{\text{RL}} = \mathcal{L}_{\text{REINFORCE}} + \lambda_{\text{CE}}\,\mathcal{L}_{\text{CE}} - \lambda_{\text{ent}}\,\mathcal{H}(\pi),
\end{equation}
where $\lambda_{\text{CE}} = 0.5$ and $\lambda_{\text{ent}} = 0.01$.
The policy gradient term uses REINFORCE with reward-to-go and an exponential moving average baseline $\bar{G}$ (decay $0.99$), with advantage normalization:
\begin{equation}
    \mathcal{L}_{\text{REINFORCE}} = -\frac{1}{T}\sum_{t=1}^{T} \log \pi_\theta(a_t \mid s_t)\,\hat{A}_t, \qquad \hat{A}_t = \frac{G_t - \bar{G}}{\sigma(G) + \epsilon},
\end{equation}
where $G_t = \sum_{\tau=t}^{T} \gamma^{\tau - t} r_\tau$ is the discounted return from step $t$ with $\gamma = 0.99$.
The auxiliary cross-entropy term $\mathcal{L}_{\text{CE}}$ is the same next-frame prediction loss as in pretraining, computed on the current observation at every RL step. This preserves the model's world-modeling capability during policy optimization.
The entropy bonus $\mathcal{H}(\pi)$ encourages exploration by penalizing overly deterministic action distributions.